\documentclass[10pt]{article}

\usepackage[utf8]{inputenc}

\usepackage[T1]{fontenc}

\usepackage{amsmath}

\usepackage{amsfonts}

\usepackage{amssymb}

\usepackage[version=4]{mhchem}

\usepackage{stmaryrd}

\usepackage{bbold}



\begin{document}

наблюдается случайный вектор $X=\left(X_{1}, X_{2}, \ldots, X_{n}\right)$, принимающий значения $x=\left(x_{1}, x_{2}, \ldots, x_{n}\right)$ в $n$-мерном евклидовом пространстве $\mathbb{R}^{n}, n \geqslant 2$, и пусть в $\mathbb{R}^{n}$ задана функция $\varphi(\cdot): \mathbb{R}^{n} \rightarrow \mathbb{R} n$, определенная по следующему правилу:



\[

\varphi(x)=x^{(\cdot)}, \quad x \in \mathbb{R}^{n},

\]



где $x^{(\cdot)}==\left(x_{(n 1)}, x_{(n 2)}, \ldots, x_{(n n)}\right)-$ вектор из $\mathbb{R}^{n}$, получен ${ }^{-}$ ный из вектора $x$ в результате перестановки его координат $x_{1}, x_{2}, \ldots, x_{n}$ в возрастающем порядке, т. е. компоненты $x_{(n 1)}, x_{(n 2)}, \ldots, x_{(n n)}$ вектора $x^{(\cdot)}$ удовлетворяют следующему соотношению



\[

x_{(n 1)} \leqslant x_{(n 2)} \leqslant \cdots \leqslant x_{(n n)} .

\]



В этом случае статистика $X^{(\cdot)}=\varphi(X)=\left(X_{(n 1)}, \ldots, X_{(n n)}\right)$ наз. в а ри а ционным р ядом (или век тором) порядко вых с т а тист ик, а ее $k$-я компонента $X_{n k}(k=1,2, \ldots, n)$ наз. $k$-й п о р я д к ов о й с т а т й с т и ко й.



В теории П. с. наиболее полно изучен случай, когда компоненты $X_{1}, X_{2}, \ldots, X_{n}$ случайного вектора $X$ суть независимые одинаково распределенные случайные величины, что в дальнейшем и будет предполагаться. Если $F(u)$ - функция распределения случайной величины $X_{i}, i=1,2, \ldots, n$, то функция распределения $F_{n k}(u) k$-й П. с. $X_{(n k)}$ вычисляется по формуле



$F_{n k}(u)=\mathrm{P}\left\{X_{(n k)} \leqslant u\right\}=I_{F(u)}(k, n-k+1),|u|<\infty,(2)$ где



\[

I_{y}(a, b)=\frac{1}{B(a, b)} \int_{0}^{y} x^{a-1}(1-x)^{b-1} d x

\]



\begin{itemize}

  \item неполная бета-функция. Из (2) следует, что если функция распределения $F(u)$ имеет плотность вероятности $f(u)$, то плотность вероятности $f_{n k}(u) k$-й П.с. $X_{(n k)}, k=1,2, \ldots, n$, тоже существует и выражается формулой

\end{itemize}



$f_{n k}(u)=\frac{n !}{(k-1) !(n-k) !}[F(u)]^{k-1}[1-F(u)]^{n-k} f(u),|u|<\infty$.



В предположении существования плотности $f(u)$ была получена совместная плотность вероятности $f_{r_{1} r_{2} \ldots r_{k}}\left(u_{1}, u_{2}, \ldots, u_{k}\right)$ П. c. $X_{\left(n r_{1}\right)}, X_{\left(n r_{2}\right)}, \ldots, X_{\left(n r_{k}\right)}$, $1 \leqslant r_{1}<r_{2}<\ldots<r_{k} \leqslant n ; k \leqslant n$, к-рая выражается формулой



\[

\begin{gathered}

f_{r_{1} r_{2} \ldots r_{k}}\left(u_{1}, u_{2}, \ldots,, u_{k}\right)=\frac{n !}{\left(r_{1}-1\right) !\left(r_{2}-r_{1}-1\right) ! \ldots\left(n-r_{k}\right) !} \times \\

\times F_{1}^{r 1}\left(u_{1}\right) f\left(u_{1}\right)\left[F\left(u_{2}\right)-F\left(u_{1}\right)\right]^{r_{2}-r_{1}-1} f\left(u_{2}\right) \ldots \\

\ldots\left[1-F\left(u_{k}\right)\right]^{n--r_{k}} f\left(u_{k}\right), \\

-\infty<u_{1}<u_{2}<\ldots<u_{k}<\infty .

\end{gathered}

\]



Формулы (2) - (4) позволяют, напр., найти распределение вероятностей т. н. әк с т р е ма льных П. с.



\[

\begin{aligned}

& X_{(n 1)}=\min _{1 \leqslant i \leqslant n}\left(X_{1}, X_{2}, \ldots, X_{n}\right) \text { и } \\

& X_{(n n)}=\max _{1 \leqslant i \leqslant n}\left(X_{1}, X_{2}, \ldots, X_{n}\right),

\end{aligned}

\]



а также распределение статистики $W_{n}=X_{(n n)}-X_{(n 1)}$, к-рую наз. р а з м а х о м. Напр., если функция распределения $F(u)$ непрерывна, то функция распределения размаха $W_{n}$ выражается формулой



$\mathrm{P}\left\{W_{n}<w\right\}=n \int_{-\infty}^{\infty}[F(u+w)-F(u)]^{n-1} d F(u), w \geqslant 0$.



Формулы (2) - (5) показывают, что, как и в общей теории выборочных методов, точные распределения П. с. невозможно использовать при получении статистич. выводов, если функция распределения $F(u)$ неизвестна. Именно поэтому в теории П. с. получили пи- рокое развитие асимптотич. методы исследования распределений П. с. при неограниченном увеличении размерности $n$ вектора наблюдений. $B$ асимптотич. теории П. с. изучаются предельные распределения соответствующим образом нормированных последовательностей П. с. $\left\{X_{(n k)}\right\}$, когда $n \rightarrow \infty ;$ при этом, вообще говоря, порядковый номер $k$ может меняться в зависимості от $n$. Если с ростом $n$ порядковый номер $k$ меняется таким образом, что существует $\lim k / n$, отличный от 0 и 1, то соответствующие П. с. $X_{(n k)}^{n \rightarrow \infty}$ рассматриваемой последовательности $\left\{X_{(n k)}\right\}$ наз. ц е н т р а л ь н ы м и или с р е д н и м и П. с. Если же $\lim _{n \rightarrow \infty} k / n$ равен 0 или 1, то П. с. $X_{(n k)}$ наз. к р а й н и м и.



В математич. статистике центральные П. с. используют при построении состоятельных последовательностей оценок для квантилей неизвестной функции распределения $F(u)$ по реализации случайного вектора $X$ или, иначе говоря, при оценивании функции $F^{-1}(u)$. Напр., пусть $x_{P}-$ квантиль уровня $P(0<P<1)$ функции распределения $F(u)$, про к-рую известно, что ее плотность вероятности $f(u)$ непрерывна и строго положительна в нек-рой окрестности точки $x_{p}$. В этом случае последовательность центральных П. с. $\left\{X_{(n k)}\right\}$ с порядковыми номерами $k=[(n+1) P+0,5]$, где $[a]-$ целая часть действительного числа $a$, является состоятельной последовательностью оценок для квантили $x_{P}, n \rightarrow \infty$. Более того, эта последовательность П. с. $\left\{X_{(n k)}\right\}$ асимптотически нормально распределена с параметрами



\[

x_{P} \text { и } \frac{P(1-P)}{f^{2}\left(x_{P}\right)(n+1)},

\]



т. е. для любого действительного числа $x$



\[

\lim _{n \rightarrow \infty} P\left\{\frac{X_{(n k)} x_{P}}{\sqrt{\frac{P(1-P)}{n+1}}} f\left(x_{P}\right)<x\right\}=\Phi(x),

\]



где Ф $(x)$ - Функция распределения стандартного нормального закона.



П р и м е р 1. Пусть $X^{(\cdot)}=\left(X_{(n 1)}, \ldots, X_{(n n)}\right)-$ вектор П. с., построенный по случайному вектору $X=$ $=\left(X_{1}, \ldots, X_{n}\right)$, компоненты к-рого суть независимые случайные величины, подчиняющиеся одному и тому же вероятностному закону, плотность вероятности к-рого непрерывна и строго положительна в нек-рой окрестности медианы $x_{1 / 2}$. В этом случае при $n \rightarrow \infty$ последовательность выборотных медиан $\left\{\mu_{n}\right\}$, определяемых для любого $n \geqslant 2$ по правилу



\[

\mu_{n}= \begin{cases}X_{(n, m+1)}, & \text { если } n \text {-нечетное qисло, } \\ \frac{1}{2}\left(X_{(n m)}+X_{(n, m+1)}\right), & \text { если } n \text {-четное число, }\end{cases}

\]



асимптотически нормально распределена с параметрами



\[

x_{1 / 2} \text { и }\left\{4(n+1) f^{2}\left(x_{1 / 2}\right)\right\}^{-1} \text {. }

\]



В частности, если



\[

f(x)=\frac{1}{V \overline{2 \pi} \sigma} \exp \left\{-\frac{(x-a)^{2}}{2 \sigma^{2}}\right\},|a|<\infty, \sigma>0,

\]



то есть $X_{i}$ подчиняется нормальному закону $N\left(a, \sigma^{2}\right)$, то в этом случае последовательность $\left\{\mu_{n}\right\}$ асимптотічески нормально распределена с параметрами $x_{1 / 2}=a$ и $\frac{\sigma^{2} \pi}{2(n+1)}$. Если последовательность статистик $\left\{\mu_{n}\right\}$ сравнить с шоследовательностью наилучших несмещенных оценок



\[

\left\{\bar{X}_{n}\right\}, \bar{X}_{n}=\frac{1}{n} \sum_{i=1}^{n} X_{i}

\]





\end{document}